Tässä diplomityössä tarkastellaan, miten Retrieval-Augmented Generation (RAG) -järjestelmiä voidaan integroida paikallisesti asennettuihin ja yrityskohtaisesti mukautettuihin tuotteen elinkaaren hallintajärjestelmiin (PLM), jotta yritysten dokumentaationhallintaan liittyviä tiedon saatavuuden haasteita voitaisiin ratkaista. Valmistavan teollisuuden organisaatioissa tieto on usein hajallaan useissa lähteissä, kuten ohjelmistomanuaaleissa, julkaisutiedotteissa ja asiakaskohtaisissa ohjeissa, mikä tekee relevantin tiedon löytämisestä hidasta ja työlästä. Samalla tiukat tiedonhallinnan ja immateriaalioikeuksien suojaan liittyvät vaatimukset rajoittavat pilvipohjaisten tekoälyratkaisujen käyttöä, sillä niitä pidetään usein riskialttiina luottamuksellisen tiedon kannalta.

Tutkimus toteutettiin Design Science Research -metodologiaa hyödyntäen ja tapaustutkimukseen nojaten. Työssä suunniteltiin, toteutettiin ja arvioitiin RAG-pohjainen tekoälyavustaja, joka integroitiiin paikallisesti asennettuun Sovelia Core -PLM-järjestelmään. Kyseinen järjestelmä on laajasti käytössä pohjoismaisissa valmistavan teollisuuden yrityksissä. Kehitetty arkkitehtuuri perustuu PostgreSQL-tietokantaan ja sen pgvector-laajennukseen vektoripohjaista hakua varten. Järjestelmä tukee useita kielimallipalveluntarjoajia abstraktiokerrosten avulla, mikä mahdollistaa joustavat käyttöönottotavat. Lisäksi ratkaisu säilyttää PLM-järjestelmän olemassa olevat käyttöoikeusmäärittelyt hyödyntämällä tietokantatason oikeuksien hallintaa. Dokumentaatioputki kokoaa ja käsittelee aineistoa useista lähteistä, kuten sisäisistä PLM-dokumenteista, toimittajan tarjoamasta dokumentaatiosta rajapinnan kautta sekä verkkopohjaisista julkaisutiedotteista. Sisältö paloiteltiin ja upotettiin automaattisesti, ja järjestelmä tukee myös versiotietoista päivitystä.

Järjestelmän arviointi yhdisti asiantuntija-arvion vastausten laadusta  sekä teknisen suorituskyvyn mittaukset kahdessa eri kielimallikokoonpanossa: pilvipohjaisessa suljetussa mallissa (OpenAI GPT-5 Nano) ja paikallisessa avoimen lähdekoodin ratkaisussa (Ollama GPT-OSS-20B erillisellä GPU-infrastruktuurilla). Asiantuntijat arvioivat 30 edustavaa kyselyä, jotka kattoivat viisi kysymystyyppiä: ominaisuudet, toimintatavat, vianetsintä, versiokohtaiset erot sekä integraatioihin liittyvät kysymykset. Vastaukset saivat keskimäärin 4,31--4,49 pistettä viisiportaisella asteikolla neljässä arviointikriteerissä: faktinen tarkkuus, relevanssi, lähteiden näkyvyys ja vastausten kattavuus. Tämä vastaa tasoa hyvän ja erinomaisen välillä. Asiantuntija-arvioiden tekoälyvastaukset generoitiin OpenAI GPT-5 Nanolla pilvipalvelussa.

Suorituskykytestit osoittivat selviä eroja vastausajoissa: pilvipohjainen ratkaisu tuotti vastaukset keskimäärin noin 26 sekunnissa, kun taas paikallinen Ollama-käyttöönotto vaati noin 91 sekuntia. Ero johtuu todennäköisesti paikallisen ratkaisun yhden GPU:n rajoitteista ja optimointieroista, eikä niinkään paikallisen käyttöönoton perusluonteisista heikkouksista. Kehittyneemmällä optimoinnilla paikallista suorituskykyä olisi mahdollista parantaa merkittävästi. Varsinaiset hakutoiminnot toimivat molemmissa kokoonpanoissa johdonmukaisesti alle 100 millisekunnissa. Laadullinen analyysi osoitti, että vastausten suurimmat puutteet johtuivat ensisijaisesti dokumentaation rajallisesta kattavuudesta, eivät niinkään haun tai generoinnin virheistä.

Tutkimuksen tulokset tarjoavat sekä teoreettisia että käytännöllisiä kontribuutioita. Teoreettisesti työ osoittaa, että dokumentaation laatu ja kattavuus ovat keskeisin yksittäinen tekijä RAG-järjestelmän toimivuudelle toimialakohtaisissa sovelluksissa. Lisäksi hybridimallit, joissa yhdistetään paikallinen infrastruktuuri ja pilvipohjainen päättely, tarjoavat realistisen tavan tasapainottaa tiedonhallinnan vaatimuksia suorituskykyä ja kustannuksia. Tapaustutkimus osoittaa myös, että yksinkertaiset hakumenetelmät (vektorien samankaltaisuus BM25-uudelleenjärjestyksellä) toimivat riittävän hyvin, kun dokumentaatio on kattavaa. Käytännön tasolla työ tuottaa tuotantovalmiin referenssitoteutuksen paikallisiin yritysympäristöihin, toimivat mallit monilähteisen dokumentaation hallintaan sekä kustannus--suorituskykyanalyysin, joka tukee käyttöönottopäätöksiä. Tulokset osoittavat, että RAG-pohjainen dokumentaatioavustin voi saavuttaa toiminnallisesti käyttökelpoisen laadun paikallisissa PLM-ympäristöissä, kun teknisen infrastruktuurin ohella panostetaan myös lähdeaineiston laatuun ja ylläpitoon.